\input _template

\setlanguage{SK}

\Title{Vtipná algebra}

\MathcompsLink{vtipna-algebra}

\Author{Patrik Bak}

\sec Úvod

V tejto zbierke úloh nájdete moje najobľúbenejšie úlohy, ktoré majú krátke zaujímavé častokrát trikové riešenia využívajúce rôzne algebraické triky. 

\sec Úlohy 

Bez ďalšieho otáľania poďme na úlohy. Kľúčom je nájsť vtipnú úpravu. Je ťažké dať všeobecný návod -- koniec-koncov, nadštandardné úlohy sú predsa len dosť o kreativite. 

\EnvId{nested-binary-operation}
\Problem{0}{}{
    Označme 
    $$
    m \circ n = \frac{mn+4}{m+n}.
    $$
    Určte
    $$
    (\dots((42 \circ 41) \circ 40) \circ \dots \circ 1) \circ 0.
    $$
}{
    V~úlohe sú zdanlivo náhodné konštanty, napríklad~$42$ vyzerá, že by mohla byť umelo nastavená. Kebyže máme ručne vyhodnocovať takto veľké čísla, zbláznime sa. Skúste si úlohu pre menšie konštanty.
}{
    Platí
    \begitems 
    \i $1 \circ 0 = 4$
    \i $(2 \circ 1) \circ 0 = 2$
    \i $((3 \circ 2) \circ 1) \circ 0 = 2$
    \i $(((4 \circ 3) \circ 2) \circ 1) \circ 0 = 2$
    \enditems
    Nie je podozrivé, že to po čase vychádza stále 2?
}{
    \Answer{$2$.}

    Všimnime si špeciálnu vlastnosť čísla $2$ vzhľadom na operáciu $\circ$. Pre ľubovoľné $m \ne -2$ platí
    $$
    m \circ 2 = \frac{2m+4}{m+2} = \frac{2(m+2)}{m+2} = 2.
    $$
    Výraz zo zadania vyhodnocujeme postupne zľava, pričom čísla vstupujú do operácie v~poradí $42, 41, \dots, 2, 1, 0$. V~momente, keď budeme medzivýsledok skladať s~číslom $2$, sa hodnota celého výrazu zmení na $2$. Všetky nasledujúce operácie budú mať tvar $2 \circ k$, čo je opäť $2$. Výsledok je preto $2$.
}

\EnvId{zero-sum-fraction-identity}
\Problem{0}{}{
    Čísla $a,b,c$ sú nenulové reálne čísla so súčtom $0$. Dokážte, že 
    $$
    \frac{1}{a+b}+\frac{1}{b+c}+\frac{1}{c+a}=\frac{a^2+b^2+c^2}{2abc}.
    $$
}{
    Podmienku $a+b+c=0$ chceme šikovne využiť, aby sme sa zbavili komplikovaných zlomkov naľavo.
}{
    Kľúčom je uvedomiť si, že $a+b=-c$, teda prvý zlomok je vlastne rovný $-\frac 1c$, podobne zvyšné. 
}{
    Po úprave stačí ukázať
    $$
    -\left(\frac 1a + \frac 1b + \frac 1c\right) = \frac{a^2+b^2+c^2}{2abc},
    $$
    po roznásobení 
    $$
    2ab+2bc+2ca=a^2+b^2+c^2.
    $$
    Teraz potrebujeme znova použiť podmienku $a+b+c=0$. Čo také by sme s~ňou mohli urobiť, aby sme dostali výrazy, o~ktorých máme niečo dokázať?
}{
    Z~podmienky $a+b+c=0$ vyplýva $a+b=-c$, $b+c=-a$ a~$c+a=-b$. Ľavú stranu dokazovanej rovnosti preto upravíme na
    $$
    -\frac{1}{c} - \frac{1}{a} - \frac{1}{b} = -\frac{ab+bc+ca}{abc}.
    $$
    Zároveň umocnením rovnosti $a+b+c=0$ na druhú dostaneme
    $$
    0 = (a+b+c)^2 = a^2+b^2+c^2 + 2(ab+bc+ca),
    $$
    z~čoho vyjadríme $ab+bc+ca = -\frac{1}{2}(a^2+b^2+c^2)$. Dosadením tohto výrazu do upravenej ľavej strany ihneď získame
    $$
    -\frac{-\frac{1}{2}(a^2+b^2+c^2)}{abc} = \frac{a^2+b^2+c^2}{2abc},
    $$
    čo sme chceli dokázať.
}

\EnvId{constrained-fraction-sum}
\Problem{0}{}{
    Čísla $a,b,c$ sú kladné reálne také, že $ab+bc+ca=1$. Určte všetky možné hodnoty výrazu 
    $$
    \frac{a(b^2+1)}{a+b}+
    \frac{b(c^2+1)}{b+c}+
    \frac{c(a^2+1)}{c+a}.
    $$
}{
    Trikom je zjednodušiť každý zlomok individuálne použitím podmienky $ab+bc+ca=1$.
}{
    Náš skúmaný výraz je plný písmen a~je tam jedno číslo, $1$. Čo takto sa ho zbaviť?
}{
    \Answer{$2$.}

    Upravme prvý zlomok dosadením $1 = ab+bc+ca$ do čitateľa:
    $$
    \gather
    \frac{a(b^2+ab+bc+ca)}{a+b} = \frac{a(b(a+b)+c(a+b))}{a+b} =\\= \frac{a(a+b)(b+c)}{a+b} = ab+ac.
    \endgather
    $$
    Cyklicky pre druhý a~tretí zlomok dostaneme $bc+ab$ a~$ca+bc$. Sčítaním všetkých troch upravených výrazov získame
    $$
    (ab+ac) + (bc+ab) + (ca+bc) = 2(ab+bc+ca).
    $$
    Keďže $ab+bc+ca=1$, hodnota výrazu je $2$.

    Ešte potrebujeme dokázať, že táto hodnota je dosiahnuteľná -- zrejme však existujú kladné čísla $a,b,c$, pre ktoré $ab+bc+ca=1$, napríklad $a=b=c=\sqrt{3}/3$.
}

\EnvId{unit-product-fraction-sum}
\Problem{0}{}{
    Čísla $a,b,c$ sú kladné reálne so súčinom 1. Určte všetky možné hodnoty výrazu 
    $$
    \frac{1}{1+a+ab}+
    \frac{1}{1+b+bc}+
    \frac{1}{1+c+ca}.
    $$
}{
    Kľúčom je vykonať nejakú šikovnú úpravu niektorých zlomkov, ktorá umožní použiť podmienku $abc=1$. Čo je taká bežná vec, ktorú vieme robiť so zlomkami?
}{
    Jednoduchý spôsob, ako využiť $abc=1$ napr. v~prvom zlomku je rozšíriť ho~$c$, potom máme
    $$
    \frac{1}{c+ac+abc}=\frac{1}{c+ac+1}.
    $$
    Tento zlomok má čírou náhodou taký istý menovateľ ako posledný zlomok. 
}{
    \Answer{$1$.}

    Prvý zlomok rozšírime číslom $c$ a~druhý výrazom $ac$. S~využitím $abc=1$ tak upravíme ich menovatele na tvar zhodný s~tretím zlomkom:
    $$
    \align
    \frac{1}{1+a+ab} &= \frac{c}{c+ac+abc} = \frac{c}{c+ac+1}, \\
    \frac{1}{1+b+bc} &= \frac{ac}{ac+abc+ab c^2} = \frac{ac}{ac+1+c}.
    \endalign
    $$
    Sčítaním všetkých troch zlomkov dostaneme
    $$
    \frac{c}{1+c+ac} + \frac{ac}{1+c+ac} + \frac{1}{1+c+ac} = \frac{c+ac+1}{1+c+ac} = 1.
    $$
    Jediná možná hodnota výrazu je $1$. Príklad $a=b=c=1$ pritom dokazuje, že táto hodnota je dosiahnuteľná.
}

\EnvId{sum-one-triple-product}
\Problem{0}{}{
    Reálne čísla $a, b, c$ majú súčet 1. Dokážte, že číslo 
    $$
    (a + bc)(b + ca)(c + ab)
    $$
    je nezáporné.
}{
    Trikom je upraviť každú zátvorku individuálne použitím podmienky $a+b+c=1$.
}{
    Dobrým indikátorom, ako upraviť každú zátvorku, je fakt, že v~jednotlivých zátvorkách máme $a+bc$, kde  je prvý člen stupňa~1 a~druhý stupňa~2 -- čo takto ich nejako oba spraviť stupňa~2?
}{
    Využijeme podmienku $a+b+c=1$ na úpravu výrazov v~zátvorke. Pre prvú zátvorku platí:
    $$
    \gather
    a + bc = a(a+b+c) + bc = a^2+ab+ac+bc =\\= a(a+b)+c(a+b) = (a+c)(a+b).
    \endgather
    $$
    Analogicky odvodíme $b+ca = (a+b)(b+c)$ a~$c+ab = (a+c)(b+c)$. Vynásobením všetkých troch výrazov dostaneme
    $$
    (a+bc)(b+ca)(c+ab) = (a+b)^2(b+c)^2(c+a)^2.
    $$
    Výsledok je súčinom druhých mocnín reálnych čísel, a~preto je nezáporný. 
}

\EnvId{reciprocal-sum-odd-powers}
\Problem{0}{}{
    Pre nenulové reálne čísla $a,b,c$ platí
    $$
    \frac1a + \frac1b + \frac1c = \frac1{a+b+c}.
    $$
    Dokážte, že pre každé nepárne prirodzené číslo $n$ platí 
    $$
    \frac1{a^n} + \frac1{b^n} + \frac1{c^n} = \frac1{a^n+b^n+c^n}.
    $$
}{
    Zabudnite na to, čo dokazujeme. Kedy vôbec môže platiť predpoklad? Skúste nájsť všetky trojice $a,b,c$, ktoré túto rovnosť spĺňajú.
}{
    Kľúčom je uvedomiť si, že po roznásobení podmienky dostaneme 
    $$
    (a+b+c)(ab+bc+ca)-abc=0.
    $$
    Tento výraz sa prekvapivo dá rozložiť na súčin. Odtiaľ to už bude priamočiare.
}{
    Upravme podmienku zo zadania vynásobením menovateľmi:
    $$
    \gather
    \frac{ab+bc+ca}{abc} = \frac{1}{a+b+c},\\
    (a+b+c)(ab+bc+ca) - abc = 0.
    \endgather
    $$
    Výraz na ľavej strane upravíme postupným vynímaním na súčin:
    $$
    \align
    (a+b+c)(ab+bc+ca) - abc &= \\ =
    (a+b+c)(ab + c(a+b)) - abc &= \\ =
    (a+b)ab + abc + (a+b)^2c + (a+b)c^2 - abc &= \\ =
    (a+b)(ab + ac + bc + c^2) &= \\ =
    (a+b)(a(b+c) + c(b+c))  &= \\ =
    (a+b)(b+c)(c+a).&
    \endalign
    $$
    Rovnosť teda platí práve vtedy, keď sú aspoň dve z~čísel opačné. Bez ujmy na všeobecnosti nech $b=-a$.
    Potom pre každé nepárne prirodzené číslo $n$ platí $b^n = (-a)^n = -a^n$.
    Ľavá strana dokazovanej rovnosti je
    $$
    \frac{1}{a^n} + \frac{1}{-a^n} + \frac{1}{c^n} = \frac{1}{c^n}.
    $$
    Pravá strana je
    $$
    \frac{1}{a^n - a^n + c^n} = \frac{1}{c^n}.
    $$
    Ľavá strana sa rovná pravej, čím je dôkaz hotový.
}

\EnvId{cyclic-reciprocal-equalities}
\Problem{0}{}{
    Navzájom rôzne reálne čísla $a,b,c$ spĺňajú
    $$
    a + \frac1b = b + \frac1c = c + \frac1a.
    $$
    Dokážte, že $|abc|=1$.
}{
    Zadanie nám hovorí, že platí
    $$
    \align 
    a + \frac 1b &= b + \frac 1c, \\
    b + \frac 1c &= c + \frac 1a, \\
    c + \frac 1a &= a + \frac 1b. \\
    \endalign
    $$
    Uvedomenie si, že ide o~sústavu troch rovníc, je dobrý krok k~tomu, aby bolo prirodzené skúšať rôzne kombinovať rovnice, ako je bežné pri práci so sústavami.
}{
    Po chvíli ľahko zistíme, že sčítanie a~odčítanie rovníc k~ničomu nevedie. Kľúčom je tieto rovnice násobiť. K~tomu však potrebujeme každú rovnicu trochu upraviť, aby sa nám pri násobení niečo vykrátilo.
}{
    Z~rovnosti $a + \frac{1}{b} = b + \frac{1}{c}$ vyjadríme rozdiel $a-b$:
    $$
    a - b = \frac{1}{c} - \frac{1}{b} = \frac{b-c}{bc}.
    $$
    Cyklicky pre ďalšie dvojice získame
    $$
    b - c = \frac{c-a}{ca} \quad\text{a}\quad c - a = \frac{a-b}{ab}.
    $$
    Keďže sú čísla $a,b,c$ navzájom rôzne, rozdiely $a-b$, $b-c$ a~$c-a$ sú nenulové. Vynásobením týchto troch rovností dostaneme
    $$
    (a-b)(b-c)(c-a) = \frac{(b-c)(c-a)(a-b)}{(abc)^2}.
    $$
    Vydelením nenulovým výrazom $(a-b)(b-c)(c-a)$ získame $1 = \frac{1}{(abc)^2}$, teda $(abc)^2 = 1$, z~čoho vyplýva $|abc|=1$.
}

\EnvId{cyclic-square-difference-system}
\Problem{0}{}{
    Pre nenulové reálne čísla $a,b,c$ platí
    $$
    a^2-b^2 = bc \hbox{\quad a \quad} b^2-c^2 = ca.
    $$
    Dokážte, že $a^2 - c^2 = ab$.
}{
    Sčítaním rovníc dostaneme presne ľavú stranu dokazovanej rovnosti. Stačí teda dokázať niečo, čo vyzerá zdanlivo jednoduchšie.
}{
    Podľa prvého návodu sme už zistili, že súčet rovníc z~podmienky má zmysel. To ešte nie je všetko, potrebujeme niečo ďalšie. Iné než sčítanie rovníc je napríklad násobenie. Predtým ich ale často musíme upraviť -- ideálne tak, aby sa po násobení niečo vykrátilo.
}{
    Sčítaním zadaných rovníc dostaneme
    $$
    a^2-c^2 = bc+ca.
    $$
    Na dôkaz tvrdenia teda stačí ukázať, že $ab = bc+ca$.
    Prvú rovnicu zo zadania upravíme na $a^2 = b(b+c)$ a~druhú na $(b-c)(b+c)=ca$. Vynásobením týchto rovností máme
    $$
    a^2(b-c)(b+c) = abc(b+c).
    $$
    Ak by $b+c=0$, tak z~prvej upravenej rovnice $a^2=0$, teda $a=0$, čo odporuje zadaniu. Preto môžeme rovnicu vydeliť nenulovým výrazom $a(b+c)$, čím dostaneme
    $$
    a(b-c) = bc,
    $$
    čo je zrejme ekvivalentné dokazovanému vzťahu $ab=bc+ca$. Tým je dôkaz hotový.
}

\EnvId{sum-of-eighth-powers}
\Problem{1}{}{
    Len s~použitím pera a~papiera (a mozgu) nájdite dve štvormiestne čísla, ktorých súčin je rovný $4^8 + 6^8 + 9^8$.
}{
    Úloha nám naznačuje, že máme hľadať rozklad na súčin. Náš výraz je ale veľmi neprehľadný, zaveďme najprv vhodnú substitúciu čísel za písmenká, aby sme ho sprehľadnili.
}{
    Vyjadrenia $4=2^2$, $9=3^2$ a $6=2 \cdot 3$ naznačujú, že kľúčová substitúcia je $a=2$, $b=3$, výraz je potom rovný 
    $$
    a^{16}+(ab)^8+b^{16}.
    $$
    Vieme toto rozložiť na súčin?
}{
    Trikom k~rozkladu $a^{16}+(ab)^8+b^{16}$ je \textit{doplnenie na štvorec}, skúste doplniť $a^{16}+b^{16}$ na štvorec. 
}{
    \Answer{Čísla $5521$ a~$8113$.}

    Nech $a=2$, $b=3$. Výraz zo zadania je potom rovný 
    $$
    a^{16} + a^8 b^8 + b^{16}.
    $$
    Tento výraz upravíme doplnením na štvorec a~následným použitím vzorca pre rozdiel štvorcov:
    $$
    \align
    a^{16}+a^8b^8+b^{16} &= (a^{16} + 2a^8b^8 + b^{16}) - a^8b^8 \\
    &= (a^8+b^8)^2 - (a^4b^4)^2 \\
    &= (a^8+b^8-a^4b^4)(a^8+b^8+a^4b^4).
    \endalign
    $$
    Dosadíme $a=2$ a~$b=3$. Potom $a^4=16$, $b^4=81$, $a^4b^4 = 1296$.
    Ďalej $a^8=256$ a~$b^8=6561$, a $a^8+b^8 = 256+6561 = 6817$.
    Jednotlivé zátvorky sú teda rovné:
    $$
    \align
    6817 - 1296 &= 5521, \\
    6817 + 1296 &= 8113.
    \endalign
    $$
    Obe nájdené čísla sú štvormiestne.
}


\EnvId{real-quartic-equation}
\Problem{1}{}{
    V~obore reálnych čísel riešte rovnicu 
    $$
    x^4+4x^3+6x^2+4x=6.
    $$
}{
    Konštanty na ľavej strane nie sú náhodné a~mali by niečo pripomínať.
}{
    To niečo, čo tieto konštanty by mali pripomínať, sa dokonca učí v~škole a má meno.
}{
    \Answer{$x=-1\pm \root 4 \of 7$.}

    Ľavá strana rovnice pripomína binomickú vetu pre 4~členy. Vskutku,
    $$
    (x+1)^4 = x^4+4x^3+6x^2+4x+1.
    $$
    To je skoro ľavá strana rovnice, potrebujeme ešte pripočítať~1. Rovnica je teda ekvivalentná s~rovnicou.
    $$
    (x+1)^4 = 7.
    $$
    V~obore reálnych čísel má rovnica dve riešenia $x=-1\pm \root 4 \of 7$.
}

\EnvId{triple-quadratic-product-equation}
\Problem{2}{}{
    V~obore reálnych čísel riešte rovnicu 
    $$
    (x^2 + 3x + 2)(x^2 - 2x - 1)(x^2 - 7x + 12) + 24 = 0.
    $$
}{
    Kvadratické členy nie sú náhodné. Čo takto sa pozrieť na každý zvlášť, nedá sa pretransformovať?
}{
    Prvým krokom úlohy je uvedomiť si, že dva z~troch našich kvadratických členov sa dajú rozložiť na súčin -- napodiv to vychádza aj pekne, je to náhoda? 
}{
    Čo so štyrmi lineárnymi členmi a~jedným kvadratickým? Trikom je z~nich znova vytvoriť tri kvadratické členy.
}{
    Aplikovaním rád z~predošlých návodov by sme mali dôjsť k~rovnici
    $$
    (x^2 - 2x - 8)(x^2 - 2x - 3)(x^2 - 2x - 1) + 24=0.
    $$
    Jej stupeň vieme zjednodušiť vhodnou substitúciou. Mala by zostať rovnica s~celočíselnými koreňmi, ktorú už vyriešime stredoškolskými postupmi. 
}{
    \Answer{Množina riešení je $\{0, 2, 1 \pm \sqrt{6}, 1 \pm 2\sqrt{2}\}$.}

    Rozložíme prvú a~tretiu zátvorku na súčin lineárnych činiteľov:
    $$
    (x+1)(x+2)(x^2-2x-1)(x-3)(x-4) + 24 = 0.
    $$
    Vhodným preskupením členov máme:
    $$
    \align
    (x+1)(x-3) &= x^2-2x-3 \\
    (x+2)(x-4) &= x^2-2x-8
    \endalign
    $$
    Dosadením do pôvodnej rovnice máme 
    $$
    (x^2-2x-3)(x^2-2x-8)(x^2-2x-1) + 24 = 0.
    $$
    Zavedieme substitúciu $y = x^2-2x$. Rovnica prejde do tvaru
    $$
    (y-3)(y-8)(y-1) + 24 = 0.
    $$
    Po roznásobení a~úprave získame
    $$
    \align
    y^3 - 12y^2 + 35y &= 0 \\
    y(y^2-12y+35) &= 0 \\
    y(y-5)(y-7) &= 0.
    \endalign
    $$
    Pre $y \in \{0, 5, 7\}$ doriešime kvadratické rovnice pre $x$:
    \begitems
    \i $x^2-2x = 0$ má korene $0, 2$.
    \i $x^2-2x-5 = 0$ má korene $1 \pm \sqrt{6}$.
    \i $x^2-2x-7 = 0$ má korene $1 \pm 2\sqrt{2}$.
    \enditems
    Množina riešení je $\{0, 2, 1 \pm \sqrt{6}, 1 \pm 2\sqrt{2}\}$.
}


\EnvId{square-differences-divisibility}
\Problem{2}{}{
    Celé čísla $a,b,c$ spĺňajú
    $$
    (a-b)^2 + (b-c)^2 + (c-a)^2 = abc.
    $$
    Dokážte, že $a+b+c+6 \mid a^3+b^3+c^3$.
}{
    Kľúčom je spomenúť si na algebraickú identitu, ktorá akosi spája niektoré z~výrazov, o~ktorých je táto úloha.
}{
    Kľúčová identita je rozklad výrazu $a^3+b^3+c^3-3abc$ na súčin. Ak to nepoznáte, skúste si tento rozklad odvodiť. V~ďalšom návode prezradíme možné triky, ako naň prísť.
}{
    K~objaveniu $a^3+b^3+c^3-3abc$ vieme použiť tieto možné postupy:
    \begitems \style a
    \i Doplníme $a^3+b^3$ na kocku, teda vyjadríme z~rozkladu $(a+b)^3$.
    \i Doplníme $a^3+b^3+c^3$ na kocku, teda vyjadríme z~rozkladu $(a+b+c)^3=...$. Z~toho extrahujeme $a^3+b^3+c^3=...$. Zostanú nám iba zmiešané členy ako $a^2b$, $ab^2$ a $abc$. Tie je potrebné vhodne usporiadať, aby sme vybrali niečo pred zátvorku.
    \enditems 
}{
    Prezradíme kľúčovú identitu
    $$
    a^3+b^3+c^3-3abc = (a+b+c)(a^2+b^2+c^2-ab-bc-ca).
    $$
    (V~predošlom návode naznačené odvodenie nájdete v~riešení.)
}{
    Ako využiť kľúčovú identitu? Potrebujeme ešte jeden trik: uvedomenie si, ako druhá zátvorka rozkladu súvisí s~podmienkou zo zadania.
}{
    Podmienku zo zadania vieme po roznásobení napísať ako 
    $$
    2(a^2+b^2+c^2-ab-bc-ca)=abc.
    $$
    Zhodou okolností sa nám tu objavila druhá zátvorka nášho rozkladu. Teraz už stačí dať všetko opatrne dohromady.
}{
    Najprv odvodíme identitu pre rozklad výrazu $a^3+b^3+c^3-3abc$. Súčet prvých dvoch kociek doplníme na kocku súčtu podľa vzťahu 
    $$
    a^3+b^3 = (a+b)^3 - 3ab(a+b).
    $$
    Následne použijeme vzorec pre súčet kociek 
    $$
    x^3+y^3=(x+y)(x^2-xy+y^2)
    $$
    na členy $(a+b)^3$ a~$c^3$:
    $$
    \align
    a^3+b^3+c^3-3abc &=\\ 
    (a+b)^3 + c^3 - 3ab(a+b) - 3abc &=\\
    (a+b+c)((a+b)^2 - (a+b)c + c^2) - 3ab(a+b+c) &=\\
    (a+b+c)(a^2+2ab+b^2 - ac - bc + c^2 - 3ab) &= \\
    (a+b+c)(a^2+b^2+c^2-ab-bc-ca)&. 
    \endalign
    $$
    Druhú zátvorku označme $K = a^2+b^2+c^2-ab-bc-ca$. Keďže $a,b,c$ sú celé čísla, aj $K$ je celé číslo.
    Podmienku zo zadania roznásobíme a~upravíme do tvaru
    $$
    (a^2-2ab+b^2) + (b^2-2bc+c^2) + (c^2-2ca+a^2) = abc,
    $$
    takže $2K=abc$. Tento vzťah dosadíme do odvodenej identity:
    $$
    a^3+b^3+c^3 - 3(2K) = (a+b+c)K.
    $$
    Osamostatnením súčtu tretích mocnín dostaneme
    $$
    a^3+b^3+c^3 = 6K + (a+b+c)K = K(a+b+c+6).
    $$
    Výraz $a+b+c+6$ teda delí $a^3+b^3+c^3$.

    \Remark Ukážeme ešte alternatívny spôsob doplnenia na kocku pomocou rozkladu ${(a+b+c)^3}$. Po roznásobení platí:
    $$
    \gather
    (a+b+c)^3 = a^3+b^3+c^3+ \\
    {}+ 3(a^2b+ab^2+b^2c+bc^2+c^2a+ca^2) + 6abc
    \endgather
    $$
    takže 
    $$
    \gather
    a^3+b^3+c^3 - 3abc = (a+b+c)^3 \\
    {}- 3(a^2b+ab^2+b^2c+bc^2+c^2a+ca^2)- 9abc,
    \endgather
    $$
    Z~posledného riadku ide vybrať~$-3$ pred zátvorku, sústreďme sa iba na vnútro tejto zátvorky:
    $$
    \align
    (a^2b+ab^2+b^2c+bc^2+c^2a+ac^2) + 3abc &= \\
    a^2b+ab^2+b^2c+c^2b+c^2a+a^2c+3abc &= \\
    [a^2b+ab^2+abc]+[b^2c+c^2b+abc]+[c^2a+a^2c+abc] &=\\
    ab(a+b+c) + bc(b+c+a) + ca(c+a+b) &= \\
    (a+b+c)(ab+bc+ca).
    \endalign
    $$
    Teraz to už len dať dokopy:
    $$
    \gather
    a^3+b^3+c^3 - 3abc = (a+b+c)^3 \\
    {}- 3(a^2b+ab^2+b^2c+bc^2+c^2a+ca^2)- 9abc = \\
    (a+b+c)^3 - 3(a+b+c)(ab+bc+ca) \\ 
    (a+b+c)((a+b+c)^2-3(ab+bc+ca)),
    \endgather
    $$
    čo je alternatívny zápis našej dokazovanej identity.
}

\DisplayHints

\DisplayProblemSolutions

\bye